\section{Hope That Transforms the Earth}

What does the certainty of heaven do to a believer walking the roads of this world? In the New Testament, fixing the mind upon heaven never breeds passivity or neglect of one's neighbor; heavenly citizenship is the furnace that forges endurance, purity, and fearless charity.

St.\ Paul sets our dual allegiance in one sentence: ``But our conversation is in heaven: from whence also we look for the Saviour, our Lord Jesus Christ'' (Phil 3:20). The Latin \textit{conversatio} (Greek \textit{politeuma}, citizenship) declares that the baptized already live under the laws of the heavenly city while sojourning abroad. To the Colossians he writes: ``Therefore if you be risen with Christ, seek the things that are above, where Christ is sitting at the right hand of God. Mind the things that are above, not the things that are upon the earth. For you are dead: and your life is hid with Christ in God. When Christ shall appear, who is your life, then you also shall appear with him in glory'' (Col 3:1--4). Against scourging, shipwreck, and daily trial, he weighs affliction on an eternal scale: ``For that which is at present momentary and light of our tribulation worketh for us above measure, exceedingly an eternal weight of glory. While we look not at the things which are seen, but at the things which are not seen. For the things which are seen are temporal: but the things which are not seen, are eternal'' (2 Cor 4:17--18). So the Epistle to the Hebrews speaks of the patriarchs: they died confessing themselves ``pilgrims and strangers on the earth,'' seeking ``a better, that is to say, a heavenly country. Therefore, God is not ashamed to be called their God: for he hath prepared for them a city'' (Heb 11:13, 16).

\subsection*{The Theological Virtue of Hope}

In the \textit{Secunda Secundae} (\textit{ST} II-II, qq.\ 17--21), St.\ Thomas Aquinas examines the supernatural virtue by which the pilgrim reaches toward that city: the theological virtue of \textbf{hope} (\textit{spes}), an infused habit of the will leaning directly on God's omnipotent mercy to attain eternal life (\textit{ST} II-II, q.\ 17, aa.\ 1, 5).\endnote{St.\ Thomas Aquinas, \textit{Summa Theologiae} II-II, q.\ 17, aa.\ 1, 2, 5; q.\ 18, a.\ 2; q.\ 20, a.\ 1; q.\ 21, a.\ 1.} Asked what hope properly seeks, Aquinas answers with quiet grandeur:

\begin{quote}
The good which we ought to hope for from God properly and chiefly is the infinite good, which is proportionate to the power of our divine helper, since it belongs to an infinite power to lead anyone to an infinite good. Such a good is eternal life, which consists in the enjoyment of God Himself. For we should hope from Him for nothing less than Himself, since His goodness, whereby He imparts good things to His creature, is no less than His Essence. Therefore the proper and principal object of hope is eternal happiness.\endnote{St.\ Thomas Aquinas, \textit{Summa Theologiae} II-II, q.\ 17, a.\ 2.}
\end{quote}

Compared with that hope, ``nothing else appears arduous'' (q.\ 17, a.\ 2, ad 3). Resting on God as both Almighty Helper and Last End, hope walks a narrow ridge between two mortal sins: \textbf{despair}, which denies God's readiness to pardon and save the penitent and becomes ``the origin of other sins'' (\textit{ST} II-II, q.\ 20, a.\ 1); and \textbf{presumption}, which despises divine justice by expecting ``to obtain pardon without repenting, or glory without merits'' (\textit{ST} II-II, q.\ 21, a.\ 1). When the pilgrim enters the Beatific Vision, hope's work ends: ``when happiness is no longer future, but present, it is incompatible with the virtue of hope. Consequently hope, like faith, is voided in heaven, and neither of them can be in the blessed'' (\textit{ST} II-II, q.\ 18, a.\ 2), though the saints still desire our salvation ``by the love of charity'' (q.\ 18, a.\ 2, ad 3).

\subsection*{The Saints' Foretaste of Heaven}

At rare moments in Christian history, God has granted His saints a fleeting foretaste of the heavenly country, shattering the illusion that this world is our home. In the autumn of 387 at the port of Ostia, five days before St.\ Monica fell into her final fever, St.\ Augustine and his mother leaned in a window overlooking an inner garden and spoke of the eternal life of the saints. Rising in thought past bodily delights, heavens, and their own minds, they touched for a heartbeat that Eternal Wisdom where nothing passes away:

\begin{quote}
We were saying, then, If to any man the tumult of the flesh were silenced,---silenced the phantasies of earth, waters, and air,---silenced, too, the poles; yea, the very soul be silenced to herself, and go beyond herself by not thinking of herself,---silenced fancies and imaginary revelations, every tongue, and every sign, and whatsoever exists by passing away \dots\ and He alone speak not by them, but by Himself, that we may hear His word, not by fleshly tongue, nor angelic voice, nor sound of thunder, nor the obscurity of a similitude, but might hear Him---Him whom in these we love---without these, like as we two now strained ourselves, and with rapid thought touched on that Eternal Wisdom which remaineth over all. If this could be sustained, and other visions of a far different kind be withdrawn, and this one ravish, and absorb, and envelope its beholder amid these inward joys, so that his life might be eternally like that one moment of knowledge which we now sighed after, were not this ``Enter thou into the joy of Thy Lord''?\endnote{St.\ Augustine, \textit{Confessions} IX.10.25, trans.\ Pilkington, NPNF1 vol.\ 1.}
\end{quote}

While they spoke, ``this world with all its delights grew contemptible to us.'' Monica asked simply, ``what do I here?'' and on her deathbed bade her sons leave her body wherever it fell: ``This only I ask, that you will remember me at the Lord's altar, wherever you be'' (\textit{Confessions} IX.10.26--11.27). Forty years later, closing \textit{The City of God} as Roman Africa crumbled under Vandal armies, Augustine looked to the eighth and eternal day consecrated by Christ's Resurrection and wrote the words Christendom has never forgotten: \textit{Ibi vacabimus et videbimus, videbimus et amabimus, amabimus et laudabimus}---``There we shall rest and see, see and love, love and praise. This is what shall be in the end without end. For what other end do we propose to ourselves than to attain to the kingdom of which there is no end?''\endnote{St.\ Augustine, \textit{The City of God} XXII.30, trans.\ Dods, collated with E.\ Hoffmann's Latin text, CSEL 40 (Vienna, 1899--1900).}

A millennium later, St.\ Teresa of Avila---Doctor of the Church, whose visions stand as approved private testimony---recorded in Chapters 38 and 40 of her \textit{Life} how brief glimpses of heaven remade her earthly walk.\endnote{St.\ Teresa of Avila, \textit{The Life of St.\ Teresa of Jesus}, ch.\ 38, \S\S 1--8, and ch.\ 40, \S\S 1--9, trans.\ David Lewis, rev.\ Benedict Zimmerman (London: Thomas Baker, 1904), pp.\ 354--357, 385--388.} Caught up while praying the Rosary, she saw her parents in heaven amid a radiance beside which ``the brightness of the sun itself seems to be something exceedingly loathsome.'' When a noblewoman in Toledo brought out gold jewels and a prized diamond to console her in illness, Teresa ``laughed to myself, and was very sorry to see what men made much of; for I thought of what our Lord had laid up for us.'' Her fear of death vanished: ``the rapture was a great help to recognise our true home, and to see that we are pilgrims here.'' In Chapter 40, shown that God is ``Truth Itself, in Itself,'' and seeing the soul as a clear mirror holding Christ at its center---clouded whenever stained by mortal sin---she rose ``with exceeding strength and earnestness of purpose to observe with all my might everything contained in the divine writings.''

St.\ Francis de Sales brings this heavenly realism into lay life in \textit{Introduction to a Devout Life} (I.16--17), teaching the believer to cry: ``Let us go, O my dear soul! let us walk towards this blessed land which is promised us: what are we doing in Egypt? I will therefore disburden myself of all such things as may divert or retard me in so happy a journey.''\endnote{St.\ Francis de Sales, \textit{Introduction to a Devout Life} I.16--17, Pustet ed., pp.\ 42--44.} Standing in imagination beside his guardian angel between heaven above and hell beneath, the Christian hears the saints call down: ``Be of good heart, dear brother; he that diligently considers the way of devotion by which we ascended hither, shall see that we acquired these immortal delights by pleasures incomparably more sweet than those of the world.''

St.\ Alphonsus Maria de Liguori adds the consolation that carries ordinary believers through sickness, obscurity, and grief. In heaven, he writes in Consideration XXIX of \textit{Preparation for Death}, the soul ``shall then understand that the tribulations, infirmities, persecutions, and losses, which she called misfortunes and divine chastisements, were all love, all means intended by Divine Providence to conduct her to Heaven.''\endnote{St.\ Alphonsus Maria de Liguori, \textit{Preparation for Death}, Consideration XXIX, Third Point (Dublin: James Duffy, 1843), pp.\ 359--361.} There every instant stays fresh: ``She shall be always satisfied, and always thirsty: for ever thirsty, and always satiated with delights: because the desire of Heaven begets no pain, and its possession produces no tediousness.'' Recalling St.\ Philip Neri tossing a cardinal's hat into the air crying ``Paradise! Paradise!'' and St.\ Mary of Egypt sustained in the desert by that single hope, Alphonsus bids every suffering Christian look up and murmur: ``Heaven! Heaven!''

\subsection*{Why Heaven Makes Charity Fearless}

Look back at the King's sentence in Matthew 25: ``Come, ye blessed of my Father, possess you the kingdom prepared for you from the foundation of the world. For I was hungry, and you gave me to eat: I was thirsty, and you gave me to drink: I was a stranger, and you took me in: naked, and you covered me: sick, and you visited me: I was in prison, and you came to me'' (Matt 25:34--36).

Whoever believes this earth is all he has must clutch his comfort and avenge every slight, for every sacrifice looks like pure loss. Whoever knows that ``our conversation is in heaven'' can give freely, forgive enemies, serve the poor, and stand firm under threat, because neither thief nor tyrant can touch his inheritance. Immediately after declaring that ``we shall see him as he is,'' St.\ John draws the moral consequence: ``And every one that hath this hope in him, sanctifieth himself, as he also is holy'' (1 John 3:3). St.\ Augustine answers directly: ``Who purifieth us but God? Yea, but God doth not purify thee if thou be unwilling. Therefore, in that thou joinest thy will to God, in that thou purifiest thyself.''\endnote{St.\ Augustine, \textit{Homilies on the First Epistle of John} IV.7, trans.\ Browne, NPNF1 vol.\ 7.}
